\hypertarget{python-3.13-free-threaded-build-gil-removal-internals}{%
\chapter{Python 3.13: Free-Threaded Build \& GIL Removal
Internals}\label{python-3.13-free-threaded-build-gil-removal-internals}}

\hypertarget{the-free-threaded-cpython-paradigm-shift}{%
\subsection{20.1 The Free-Threaded CPython Paradigm
Shift}\label{the-free-threaded-cpython-paradigm-shift}}

Introduced provisionally in Python 3.13 (PEP 703), the free-threaded
build of CPython (also known as the ``no-GIL'' build) allows executing
bytecode in parallel across multiple OS threads without the constraint
of the Global Interpreter Lock (GIL).

In standard CPython, the GIL serves as a single global mutex protecting
the interpreter state and all Python objects from concurrent access.
This guarantees that only one thread executes bytecode at a time,
simplifying the VM design and preventing memory races. However, it
severely limits Python's ability to scale on multi-core processors.

To remove the GIL safely, the CPython runtime has been fundamentally
re-engineered. Moving from single-threaded assumptions to a scalable
parallel execution engine required redesigning three pillars of the
runtime: 1. \textbf{Reference Counting}: Overcoming atomic bus
contention on object refcount updates. 2. \textbf{Memory Allocation}:
Ensuring lock-free, thread-local allocations for object instantiation.
3. \textbf{Collection Mutability}: Protecting built-in data structures
(dictionaries, lists, sets) from race conditions.

\begin{center}\rule{0.5\linewidth}{0.5pt}\end{center}

\hypertarget{biased-reference-counting-brc-and-object-headers}{%
\subsection{20.2 Biased Reference Counting (BRC) and Object
Headers}\label{biased-reference-counting-brc-and-object-headers}}

Reference counting is the primary mechanism for memory management in
CPython. In a standard build, incrementing/decrementing reference counts
is a simple, non-atomic operation:
\passthrough{\lstinline!op->ob\_refcnt++!}. In a thread-parallel
environment, standard operations would cause race conditions. However,
using standard atomic operations (\passthrough{\lstinline!atomic\_add!}
/ \passthrough{\lstinline!atomic\_sub!} or
\passthrough{\lstinline!LOCK XADD!} instructions) globally would destroy
performance due to cache coherency traffic (bus locks) across CPU cores.

To solve this, PEP 703 introduces \textbf{Biased Reference Counting
(BRC)}. BRC divides an object's reference count state based on thread
ownership.

\hypertarget{biased-object-headers}{%
\subsubsection{1. Biased Object Headers}\label{biased-object-headers}}

An object is ``biased'' toward the thread that allocated it (the owner
thread). The owner thread uses fast, non-atomic CPU instructions to
modify the reference count. Non-owner threads (foreign threads) must use
atomic operations on a separate ``shared'' reference count field.

Under no-GIL configurations, the standard
\passthrough{\lstinline!PyObject!} header
(\passthrough{\lstinline!Include/object.h!}) is redefined:

\begin{lstlisting}[language=C]
typedef struct _object {
    uintptr_t ob_ref_local;     /* Local reference count and owner thread ID */
    uintptr_t ob_ref_shared;    /* Shared reference count and object status flags */
    PyTypeObject *ob_type;      /* Pointer to object type descriptor */
} PyObject;
\end{lstlisting}

\hypertarget{bit-layout-and-fields-representation}{%
\subsubsection{2. Bit Layout and Fields
representation}\label{bit-layout-and-fields-representation}}

The fields \passthrough{\lstinline!ob\_ref\_local!} and
\passthrough{\lstinline!ob\_ref\_shared!} encode multiple pieces of
metadata to optimize memory usage:

\begin{itemize}
\tightlist
\item
  \passthrough{\lstinline!ob\_ref\_local!}:

  \begin{itemize}
  \tightlist
  \item
    \textbf{Thread ID Bits (Upper Bits)}: Holds the memory address of
    the owner thread's thread-state structure
    (\passthrough{\lstinline!PyThreadState*!}).
  \item
    \textbf{Local Refcount (Lower Bits)}: A small counter tracking
    reference updates made by the owner thread. Typically, the lower 16
    or 32 bits are used for this count.
  \end{itemize}
\item
  \passthrough{\lstinline!ob\_ref\_shared!}:

  \begin{itemize}
  \tightlist
  \item
    \textbf{Shared Refcount (Upper Bits)}: A signed counter tracking
    reference additions/subtractions made by non-owner threads.
  \item
    \textbf{Status Flags (Lower Bits)}: Bits reserved for encoding
    object states:

    \begin{itemize}
    \tightlist
    \item
      \passthrough{\lstinline!STATE\_STATIC!} (bit 0): Set if the object
      is static (e.g., statically allocated builtin types).
    \item
      \passthrough{\lstinline!STATE\_IMMORTAL!} (bit 1): Set if the
      object is immortal (reference counts are never modified).
    \item
      \passthrough{\lstinline!STATE\_DEFERRED!} (bit 2): Set if the
      object uses deferred reference counting (GC-managed objects).
    \end{itemize}
  \end{itemize}
\end{itemize}

\hypertarget{brc-reference-counting-algorithm}{%
\subsubsection{3. BRC Reference Counting
Algorithm}\label{brc-reference-counting-algorithm}}

When modifying a reference (via \passthrough{\lstinline!Py\_INCREF!} or
\passthrough{\lstinline!Py\_DECREF!}), the runtime performs owner-thread
detection:

\[\text{Owner Thread State Pointer} = \text{ob\_ref\_local} \& \text{THREAD\_STATE\_MASK}\]

\begin{lstlisting}
                   [Reference Count Change Request]
                                  |
               Get Current ThreadState (Tstate)
                                  |
              Is Tstate == (ob_ref_local & MASK)?
                                /   \
                              Yes    No
                              /       \
      [Non-Atomic Local Increment]   [Atomic Shared Increment]
       (Fast path: CPU registers)     (Slow path: CAS / Bus Lock)
\end{lstlisting}

\hypertarget{increment-code-walkthrough}{%
\paragraph{Increment Code
Walkthrough}\label{increment-code-walkthrough}}

\begin{lstlisting}[language=C]
static inline void
_Py_INCREF_Specialized(PyObject *op, PyThreadState *tstate)
{
    uintptr_t local = op->ob_ref_local;
    // Fast path: current thread is the owner
    if ((local & _Py_THREAD_ID_MASK) == (uintptr_t)tstate) {
        op->ob_ref_local = local + _Py_REF_INCREMENT;
    }
    // Slow path: foreign thread
    else {
        _Py_Incref_Shared(op);
    }
}
\end{lstlisting}

\hypertarget{decrement-and-deallocation-walkthrough}{%
\paragraph{Decrement and Deallocation
Walkthrough}\label{decrement-and-deallocation-walkthrough}}

When a reference is decremented, if the local count reaches zero, the
owner thread merges the shared reference count:

\begin{lstlisting}[language=C]
static inline void
_Py_DECREF_Specialized(PyObject *op, PyThreadState *tstate)
{
    uintptr_t local = op->ob_ref_local;
    if ((local & _Py_THREAD_ID_MASK) == (uintptr_t)tstate) {
        uintptr_t new_local = local - _Py_REF_INCREMENT;
        if ((new_local & _Py_REF_MASK) == 0) {
            // Local count hit zero; merge shared reference counts to check for deallocation
            _Py_Merge_And_Dealloc(op);
        } else {
            op->ob_ref_local = new_local;
        }
    } else {
        _Py_Decref_Shared(op);
    }
}
\end{lstlisting}

During \passthrough{\lstinline!\_Py\_Merge\_And\_Dealloc(op)!}, the
owner thread atomically reads and clears
\passthrough{\lstinline!ob\_ref\_shared!}. If the sum of
\passthrough{\lstinline!ob\_ref\_local!} (cleared of thread state bits)
and \passthrough{\lstinline!ob\_ref\_shared!} is less than or equal to
zero, it calls the type's \passthrough{\lstinline!tp\_dealloc!} slot to
free the object.

\begin{center}\rule{0.5\linewidth}{0.5pt}\end{center}

\hypertarget{thread-safe-allocation-via-mimalloc}{%
\subsection{20.3 Thread-Safe Allocation via
mimalloc}\label{thread-safe-allocation-via-mimalloc}}

CPython's classic allocator (\passthrough{\lstinline!pymalloc!}) is
optimized for small objects but relies on single-threaded assumptions.
Under a free-threaded runtime, sharing a global allocator lock would
lead to lock contention. Consequently, PEP 703 replaces
\passthrough{\lstinline!pymalloc!} with \textbf{mimalloc}, a
thread-safe, metadata-compact allocator developed by Microsoft.

\hypertarget{mimalloc-structural-hierarchy}{%
\subsubsection{1. Mimalloc Structural
Hierarchy}\label{mimalloc-structural-hierarchy}}

Mimalloc organizes memory allocations into hierarchical structures to
minimize thread synchronization:

\begin{lstlisting}
+-------------------------------------------------------+
|                       OS Page                         |
|  +-------------------------------------------------+  |
|  |                    Segment                      |  |
|  |  +------------------+     +------------------+  |  |
|  |  |      Page 1      |     |      Page 2      |  |  |
|  |  |  +------------+  |     |  +------------+  |  |  |
|  |  |  | Local Heap |  |     |  | Local Heap |  |  |  |
|  |  |  +------------+  |     |  +------------+  |  |  |
|  |  +------------------+     +------------------+  |  |
|  +-------------------------------------------------+  |
+-------------------------------------------------------+
\end{lstlisting}

\begin{itemize}
\tightlist
\item
  \textbf{Segments (typically 4MB)}: Large contiguous memory blocks
  requested from the operating system.
\item
  \textbf{Pages (typically 64KB - 512KB)}: Subdivisions of segments
  containing blocks of a single size class (e.g., 32-byte blocks,
  64-byte blocks).
\item
  \textbf{Heaps}: Thread-local handles containing lists of active pages.
\end{itemize}

\hypertarget{thread-local-allocations}{%
\subsubsection{2. Thread-Local
Allocations}\label{thread-local-allocations}}

Every OS thread maintains its own thread-local mimalloc heap
(\passthrough{\lstinline!mi\_heap\_t!}). When a thread executes
\passthrough{\lstinline!PyObject\_New()!}, it attempts to allocate from
its thread-local heap: 1. It locates the active page corresponding to
the requested size class. 2. It pops a block from the page's
thread-local \textbf{free list}. This list is accessed only by the
owning thread, requiring no atomic locks or memory barriers.

\hypertarget{cross-thread-deallocations-and-thread-safe-free-lists}{%
\subsubsection{3. Cross-Thread Deallocations and Thread-Safe Free
Lists}\label{cross-thread-deallocations-and-thread-safe-free-lists}}

When a foreign thread deallocates an object, writing directly back to
the owner thread's thread-local free list would cause race conditions.
To resolve this, mimalloc uses a secondary, atomic free list for each
page:

\begin{lstlisting}
[Foreign Thread deallocates block]
               |
               v (Atomic CAS operation)
      +------------------+
      |  Atomic Free List| (Attached to Page)
      +------------------+
               |
               v (During thread-local heap maintenance)
      +------------------+
      | Local Free List  | (Accessed non-atomically by Owner Thread)
      +------------------+
\end{lstlisting}

When a foreign thread frees memory: 1. It atomically prepends the freed
block to the page's \textbf{atomic free list}
(\passthrough{\lstinline!thread\_free!}) using a lock-free
Compare-And-Swap (CAS) loop. 2. When the owner thread exhausts its
thread-local \passthrough{\lstinline!free!} list, it cleans up and
merges the \passthrough{\lstinline!thread\_free!} list into its local
list. This deferred merging drastically reduces cache line bouncing
across CPU cores.

\begin{center}\rule{0.5\linewidth}{0.5pt}\end{center}

\hypertarget{lock-free-collections-thread-safety}{%
\subsection{20.4 Lock-free Collections \& Thread
Safety}\label{lock-free-collections-thread-safety}}

In standard CPython, code executing modifications to dictionaries,
lists, or sets does not need to worry about internal consistency, as the
GIL serializes all operations. Without the GIL, concurrent reads/writes
could corrupt collection structures, leading to segmentation faults or
memory corruption.

\hypertarget{dictionary-locking-system}{%
\subsubsection{1. Dictionary Locking
System}\label{dictionary-locking-system}}

CPython's compact dictionary layout was modified to support concurrent
read and write operations. * \textbf{Read Paths (Lock-free)}: Operations
like dict lookups (\passthrough{\lstinline!PyDict\_GetItem!}) read the
hash index and table entries without acquiring locks. They rely on
atomic pointer reads and memory barriers to ensure they read consistent
states. * \textbf{Write Paths (Fine-grained Locks)}: Modifying
operations (insertion, deletion) acquire a localized lock associated
with the dictionary instance. Instead of locking the global runtime,
CPython locks only the specific dictionary being modified.

To avoid allocating a full OS mutex for every dictionary (which would
consume excessive memory), CPython uses a pool of shared locks, indexing
into them using a hash of the dictionary's memory address.

\hypertarget{lock-arrays-and-pymutex}{%
\subsubsection{2. Lock Arrays and
PyMutex}\label{lock-arrays-and-pymutex}}

To support light, low-overhead synchronization, CPython implements a
highly efficient locking primitive: \textbf{PyMutex}.

\begin{lstlisting}[language=C]
typedef struct {
    uint8_t v; /* Lock state byte */
} PyMutex;
\end{lstlisting}

\begin{itemize}
\tightlist
\item
  \textbf{State Byte \passthrough{\lstinline!v!}}:

  \begin{itemize}
  \tightlist
  \item
    Bit 0: Indicates if the lock is held (1) or free (0).
  \item
    Bit 1: Indicates if there are threads waiting in the queue (1) or
    not (0).
  \end{itemize}
\item
  \textbf{Fast Path}: A thread attempts to acquire the lock using a
  single atomic test-and-set instruction
  (\passthrough{\lstinline!atomic\_compare\_exchange!} on bit 0). If
  successful, it proceeds without sleeping or system calls.
\item
  \textbf{Slow Path}: If the lock is held, the thread registers itself
  in a global wait queue associated with the lock's memory address and
  suspends execution, waiting to be woken up by the lock holder.
\end{itemize}

\begin{center}\rule{0.5\linewidth}{0.5pt}\end{center}

\hypertarget{generational-garbage-collection-gc-in-a-gil-less-world}{%
\subsection{20.5 Generational Garbage Collection (GC) in a GIL-less
World}\label{generational-garbage-collection-gc-in-a-gil-less-world}}

CPython's cyclic garbage collector identifies self-referencing loops
that reference counting alone cannot reclaim. In a standard build, the
GC runs synchronously on the main thread during allocations. In a
free-threaded environment, running the GC requires coordinated
synchronization across all running threads to prevent mutator threads
from modifying object references during the collection pass.

\hypertarget{stop-the-world-stw-synchronization}{%
\subsubsection{1. Stop-the-World (STW)
Synchronization}\label{stop-the-world-stw-synchronization}}

To run the cyclic collection safely, CPython implements a Stop-the-World
mechanism:

\begin{lstlisting}
[GC Thread initiates GC] ---> Sets global GC status flag
                                      |
     [All other threads reach safe-points / bytecode loop boundary]
                                      |
                          Suspends all mutator threads
                                      |
                          Executes GC Collection Pass
                                      |
                           Resumes all mutator threads
\end{lstlisting}

\begin{enumerate}
\def\labelenumi{\arabic{enumi}.}
\tightlist
\item
  A thread triggering a collection sets a global GC request flag.
\item
  All other running threads check this flag at defined bytecode
  boundaries (safe-points).
\item
  Upon detecting the flag, threads suspend their execution and block on
  an OS condition variable.
\item
  Once all threads are verified as stopped, the GC thread safely
  performs the collection pass, scanning object references and breaking
  cyclic loops.
\item
  After completion, the GC thread signals the condition variable,
  resuming all suspended threads.
\end{enumerate}

\hypertarget{deferred-reference-counting-drc}{%
\subsubsection{2. Deferred Reference Counting
(DRC)}\label{deferred-reference-counting-drc}}

For objects created during compilation (such as functions, classes, code
objects, and constants), standard reference counting is completely
bypassed. These objects are marked as
\passthrough{\lstinline!STATE\_DEFERRED!} or
\passthrough{\lstinline!STATE\_IMMORTAL!} inside
\passthrough{\lstinline!ob\_ref\_shared!}. * \textbf{Immortal Objects}:
Reference count updates are completely skipped. They are never
deallocated during the lifetime of the process. * \textbf{Deferred
Objects}: Reference count updates are ignored at runtime. The GC takes
full responsibility for tracking these objects and determining when they
are no longer reachable, freeing their memory during GC cycles. This
bypasses atomic reference operations on highly accessed structural
components.

\begin{center}\rule{0.5\linewidth}{0.5pt}\end{center}
